% id: 94-32
% book: 베이직쎈 대수 · 05 지수함수와 로그함수의 활용
% section: 유형 057 로그부등식의 실생활에의 활용
% figure: none
% answer: 
% answer_source: 
% variant_level: 0 (원본 전사)
% 이 파일은 build.mjs 가 items.json 에서 만든다. 고칠 때는 items.json 을 고친다.
\smash{\raisebox{1.5mm}{\dmkichul{베이직쎈 대수 94쪽 32번}}}
화재가 발생한 실내의 온도는 시간에 따라 변한다. 초기 온도가 $15\mkern3mu ^{\circ}\mathrm{C}$인 어떤 실내에서 화재가 발생한 지 $t$분 후 온도를 $T\mkern3mu ^{\circ}\mathrm{C}$라 하면 다음과 같은 관계식이 성립한다고 한다.\par\smallskip\dispeq{$T=15+345\log (8t+1)$}\par\smallskip\noindent 이 실내에서 화재가 발생한 지 $a$분 후 처음으로 온도가 $360\mkern3mu ^{\circ}\mathrm{C}$ 이상이었을 때, $a$의 값은?
\dmchoicesauto{\rule[-3ex]{0pt}{8ex}}{$1$}{$\dfrac{9}{8}$}{$\dfrac{5}{4}$}{$\dfrac{11}{8}$}{$\dfrac{3}{2}$}
